% META: {"id": "1SPE-SUITES-CO-028", "chapitre": "1SPE-SUITES", "type_objet": "corrige", "exercice_ref": "1SPE-SUITES-EX-028", "capacites_codes": ["C2", "C5"], "sources_inspiration": [], "status": "generated"}
% BEGIN-VERIFY
% from sympy import symbols, simplify
% n = symbols('n', integer=True, nonnegative=True)
% u = -4*n + 30
% assert u.subs(n, 0) == 30
% assert u.subs(n, 1) == 26
% assert u.subs(n, 2) == 22
% diff = u.subs(n, n+1) - u
% assert simplify(diff) == -4
% # Suite decroissante car r = -4 < 0
% # u_n > 0 => -4n + 30 > 0 => n < 7.5 => n <= 7
% assert u.subs(n, 7) == 2
% assert u.subs(n, 8) == -2
% # v_n = u_n^2
% v = u**2
% assert v.subs(n, 0) == 900
% assert v.subs(n, 1) == 676
% assert v.subs(n, 2) == 484
% # v_0 > v_1 > v_2 : décroissant au debut, mais v_7 = 4, v_8 = 4, v_9 = 36
% assert v.subs(n, 7) == 4
% assert v.subs(n, 8) == 4
% assert v.subs(n, 9) == 36
% END-VERIFY
\begin{corrige}{1SPE-SUITES-EX-028}

\textbf{Question 1 — Calcul de $u_0$, $u_1$ et $u_2$.}

\[
  u_0 = -4 \times 0 + 30 = 30, \quad u_1 = -4 \times 1 + 30 = 26, \quad u_2 = -4 \times 2 + 30 = 22.
\]

\medskip
\textbf{Question 2 — Démonstration que $(u_n)$ est arithmétique.}

Pour tout entier naturel $n$, on calcule la différence :
\[
  u_{n+1} - u_n = \bigl(-4(n+1) + 30\bigr) - (-4n + 30) = -4n - 4 + 30 + 4n - 30 = -4.
\]

\commentaireMarge{La différence $u_{n+1} - u_n$ est constante : c'est la définition d'une suite arithmétique.}

La différence $u_{n+1} - u_n = -4$ est constante pour tout $n \in \mathbb{N}$.

Donc $(u_n)$ est une suite arithmétique de raison $r = -4$ et de premier terme $u_0 = 30$.

\medskip
\textbf{Question 3 — Sens de variation.}

La raison de $(u_n)$ est $r = -4 < 0$.

\commentaireMarge{Une suite arithmétique est strictement décroissante si et seulement si sa raison est strictement négative.}

Puisque $r < 0$, pour tout $n \in \mathbb{N}$ on a $u_{n+1} - u_n = -4 < 0$, donc $u_{n+1} < u_n$.

La suite $(u_n)$ est \textbf{strictement décroissante}.

\medskip
\textbf{Question 4 — Valeurs de $n$ telles que $u_n > 0$ et $u_n \leq 0$.}

On résout $u_n > 0$ :
\[
  -4n + 30 > 0 \iff -4n > -30 \iff n < \frac{30}{4} = 7{,}5.
\]

Comme $n$ est un entier naturel, $u_n > 0 \iff n \leq 7$.

\textit{Vérification :} $u_7 = -28 + 30 = 2 > 0$ et $u_8 = -32 + 30 = -2 < 0$. \checkmark

De même, $u_n \leq 0 \iff n \geq 8$.

\commentaireMarge{$u_n > 0$ pour $n \in \{0, 1, \ldots, 7\}$ et $u_n \leq 0$ pour tout entier $n \geq 8$.}

\medskip
\textbf{Question 5 — Étude de $(v_n)$ avec $v_n = u_n^2$.}

On calcule les premiers termes :
\[
  v_0 = 30^2 = 900, \quad v_1 = 26^2 = 676, \quad v_2 = 22^2 = 484.
\]

Ces trois valeurs sont dans l'ordre décroissant ($v_0 > v_1 > v_2$).

Cependant, $(v_n)$ \textbf{n'est pas monotone} : en effet, $u_n$ devient négatif à partir de $n = 8$, et $|u_n|$ commence à croître. Par exemple :
\[
  v_7 = 2^2 = 4, \quad v_8 = (-2)^2 = 4, \quad v_9 = (-6)^2 = 36.
\]

\commentaireMarge{On observe que $v_7 = v_8 = 4$ puis $v_9 = 36 > v_8$ : la suite cesse d'être décroissante. Le carré d'une fonction décroissante n'est pas monotone en général.}

La suite $(v_n)$ est décroissante tant que $u_n$ est positive (pour $n \leq 7$), puis elle reprend une croissance lorsque $|u_n|$ augmente (pour $n \geq 8$). Elle n'est donc pas monotone.

\baremeIndicatif{Q1 : 1 pt (calculer : trois termes) ; Q2 : 2 pts (raisonner : calcul de la différence 1 pt + conclusion 1 pt) ; Q3 : 1 pt (raisonner : signe de $r$) ; Q4 : 2 pts (calculer, raisonner : résolution de l'inéquation 1 pt + vérification 1 pt) ; Q5 : 2 pts (raisonner, communiquer : calculs de $v_0$, $v_1$, $v_2$ 1 pt + explication de la non-monotonicité 1 pt)}
\end{corrige}
